
La creación de circuitos biohíbridos tienen diversas aplicaciones tecnológicas. Entre ellas destacan la \textit{Organoid Intelligence}, disciplina en la cual se usan células neuronales para procesar información, usando una técnica similar al \textit{machine learning}, pero aplicada a neuronas reales, lo que facilita la resolución de problemas complejos \cite{Smirnova2023}. Creación de dispositivos \textit{SLC-MLC}, que hacen partícipes a las propias células en el almacenamiento de datos, mejorando en algunos casos el rendimiento y el coste energético \cite{Murugan2012}. Además ayudan al estudio y la comprensión de las neuronas y las estructuras funcionales que estas forman, facilitando en un futuro la creación de dispositivos que ayuden a personas con condiciones neurológicas a controlar o paliar con las mismas.
